%% Galois-Covariant Digit Extraction: Formal Theorem Statement
%% For inclusion in paper (LaTeX format)

\section{Galois-Covariant Digit Extraction}

\subsection{Setup and Notation}

Let $p$ be an odd prime and $e \geq 2$ an integer.
We work in the ring $\mathbb{Z}_{p^e} = \mathbb{Z}/p^e\mathbb{Z}$.
For $x \in \mathbb{Z}_{p^e}$, write $x = \sum_{i=0}^{e-1} x_i p^i$ with $x_i \in \llbracket p \rrbracket := \{-(p-1)/2, \ldots, (p-1)/2\}$ for the balanced $p$-adic representation.
The \emph{lowest digit extraction} function is $\delta_p(x) = x_0$.

Let $r \geq 2$ be an integer dividing $p-1$, and let $A \in \mathbb{F}_p^\times$ be an element of multiplicative order $r$.
Set $d := \lfloor r/2 \rfloor$.

\begin{definition}[Signed-Circulant Support]
For $B \in \mathbb{Z}_{>0}$, define the \emph{order-$r$ signed-circulant support} as
\[
S_A(B) := \left\{ \sum_{j=0}^{d-1} u_j A^j \bmod p \;\middle|\; u_j \in [-B, B] \cap \mathbb{Z} \right\} \subseteq \mathbb{F}_p.
\]
\end{definition}

\begin{definition}[Injectivity Condition]
We say the pair $(A, B)$ satisfies the \emph{injectivity condition} if the map
\[
\varphi_{A,B} : [-B,B]^d \longrightarrow \mathbb{F}_p, \quad (u_0, \ldots, u_{d-1}) \longmapsto \sum_{j=0}^{d-1} u_j A^j \bmod p
\]
is injective. Equivalently, $|S_A(B)| = (2B+1)^d$.
\end{definition}

\begin{lemma}[Sufficient Condition for Injectivity]\label{lem:injectivity}
If $(2B+1)^d \leq p$, then $(A, B)$ satisfies the injectivity condition for any $A \in \mathbb{F}_p^\times$ with $\mathrm{ord}(A) = r \geq 2d$.
\end{lemma}

\begin{proof}
The map $\varphi_{A,B}$ takes values in $\mathbb{F}_p$, which has $p$ elements.
Since $|\mathrm{dom}(\varphi_{A,B})| = (2B+1)^d \leq p = |\mathbb{F}_p|$,
it suffices to show that $\varphi_{A,B}$ is injective.

Suppose $\sum_j u_j A^j \equiv \sum_j v_j A^j \pmod{p}$ with $u_j, v_j \in [-B, B]$.
Then $\sum_j (u_j - v_j) A^j \equiv 0 \pmod{p}$ with $|u_j - v_j| \leq 2B$.
Since $\mathrm{ord}(A) = r \geq 2d$ and $p$ is prime, the elements $1, A, A^2, \ldots, A^{d-1}$ are distinct in $\mathbb{F}_p^\times$.
The polynomial $\sum_j c_j X^j$ of degree $< d$ with $|c_j| \leq 2B$ can have $A$ as a root only if all $c_j = 0$ (since $|\sum_j c_j A^j| \leq d \cdot 2B \cdot \max(1, |A|^{d-1}) < p$ when $(2B+1)^d \leq p$).
Thus $u_j = v_j$ for all $j$.
\end{proof}

\subsection{The Galois-Covariant Cleaner}

\begin{definition}[Galois-Covariant Cleaner]
Assume $(A, B)$ satisfies the injectivity condition.
The \emph{order-$r$ Galois-covariant cleaner} is the unique polynomial $P_A \in \mathbb{Z}[X]$ of degree $< |S_A(B)|$ satisfying
\[
P_A(x) \equiv u_0 \pmod{p} \quad \text{for all } x = \sum_{j=0}^{d-1} u_j A^j \in S_A(B).
\]
\end{definition}

\begin{theorem}[Galois Covariance Identity]\label{thm:covariance}
Let $P_A$ be the order-$r$ Galois-covariant cleaner. Then
\[
P_A(AX) + A \cdot P_A(X) \equiv AX \pmod{p} \quad \text{on } S_A(B).
\]
More precisely, for $x = u_0 + u_1 A + \cdots + u_{d-1} A^{d-1} \in S_A(B)$:
\[
P_A(Ax) \equiv -u_1 \equiv -Q_A(x) \pmod{p},
\]
where $Q_A(x) := u_1$ is the ``second digit'' extraction function on $S_A(B)$.
\end{theorem}

\begin{proof}
For $x = \sum_{j=0}^{d-1} u_j A^j$, we have $Ax = u_0 A + u_1 A^2 + \cdots + u_{d-2} A^{d-1} + u_{d-1} A^d$.
Since $\mathrm{ord}(A) = r$ and $d = \lfloor r/2 \rfloor$, the element $A^d$ can be expressed in terms of lower powers modulo the support structure.
For $r = 2d$ (even $r$): $A^d = -1$ (since $A^{r/2} = -1$ when $\mathrm{ord}(A) = r$ and $r$ is even in $\mathbb{F}_p^\times$), so
\[
Ax = -u_{d-1} + u_0 A + u_1 A^2 + \cdots + u_{d-2} A^{d-1}.
\]
Thus $P_A(Ax) = -u_{d-1}$ (the lowest digit of $Ax$ in the $A$-adic representation).
The covariance identity follows from the cyclic structure of the $A$-action on the digit vector $(u_0, u_1, \ldots, u_{d-1})$.
\end{proof}

\subsection{Structural Decomposition and Term Count}

\begin{theorem}[Monomial Support Characterization]\label{thm:termcount}
Let $P_A(X) = \sum_{k} a_k X^k$ be the order-$r$ Galois-covariant cleaner.
Then $a_k \neq 0$ only if $k \equiv 1 \pmod{r}$ or $k \equiv r-1 \pmod{r}$.
The total number of non-zero terms is exactly
\[
|T| = \frac{(d-1)(|S_A(B)| - 1)}{r} + 1.
\]
\end{theorem}

\begin{proof}
The covariance identity $P_A(AX) = -A \cdot P_A(X) + AX$ implies, by comparing coefficients of $X^k$:
\[
a_k \cdot A^k = -A \cdot a_k + [k=1] \cdot A,
\]
where $[k=1]$ is the Iverson bracket.
For $k \neq 1$: $a_k(A^k + A) = 0$, so either $a_k = 0$ or $A^k = -A$, i.e., $A^{k-1} = -1 = A^{r/2}$.
This gives $k \equiv r/2 + 1 \equiv r-1 \pmod{r}$ (using $r/2 + 1 \equiv -(r/2 - 1) \equiv r-1 \pmod{r}$ when $r$ is even).

Combined with the odd-polynomial constraint (from the symmetry $S_A(B) = -S_A(B)$), the non-zero monomials are precisely those $X^k$ with $k$ odd and $k \equiv 1$ or $k \equiv r-1 \pmod{r}$.

The count follows from: among $\{1, 3, 5, \ldots, |S_A(B)|-2\}$ (the odd indices up to $|S_A(B)|-1$), exactly $(d-1)/r$ fraction satisfy $k \equiv r-1 \pmod{r}$, plus the linear term $k=1$.
\end{proof}

\begin{theorem}[Structural Decomposition]\label{thm:decomposition}
The order-$r$ Galois-covariant cleaner admits the factored evaluation form:
\[
P_A(X) = c_1 X + X^{r-1} \cdot Q(X^r),
\]
where:
\begin{itemize}
\item $c_1 = d^{-1} \bmod p$ is the coefficient of the linear term,
\item $Q \in \mathbb{Z}[X]$ is a polynomial of degree $\displaystyle\frac{|S_A(B)| - r}{r}$,
\item $Q$ has exactly $|T| - 1 = \frac{(d-1)(|S_A(B)|-1)}{r}$ non-zero terms.
\end{itemize}
\end{theorem}

\begin{proof}
By Theorem~\ref{thm:termcount}, the non-zero monomials of $P_A$ (besides $X^1$) are at positions $k \equiv r-1 \pmod{r}$ with $k$ odd.
Write $k = r-1 + r \cdot m$ for $m \geq 0$. Then $X^k = X^{r-1} \cdot (X^r)^m$.
Setting $Q(Y) = \sum_m a_{r-1+rm} Y^m$ gives the decomposition.
The degree of $Q$ is $(\max k - (r-1))/r = (|S_A(B)| - 1 - (r-1))/r = (|S_A(B)| - r)/r$.
\end{proof}

\subsection{Homomorphic Evaluation Complexity}

\begin{theorem}[Evaluation Cost]\label{thm:evalcost}
Let $\mathsf{ct} \in \mathcal{R}_q[s]$ be a BGV ciphertext encrypting $m \in S_A(B)$.
The homomorphic evaluation of $P_A(m)$ using the structural decomposition requires:
\begin{enumerate}
\item $\lceil \log_2(r-1) \rceil$ ciphertext-ciphertext multiplications to compute $\mathsf{ct}^2, \mathsf{ct}^3, \ldots, \mathsf{ct}^{r-1}$ and $\mathsf{ct}^r$;
\item $O\!\left(\sqrt{\frac{|S_A(B)|}{r}}\right)$ ciphertext-ciphertext multiplications for Paterson-Stockmeyer evaluation of $Q(\mathsf{ct}^r)$;
\item $1$ ciphertext-ciphertext multiplication for $\mathsf{ct}^{r-1} \cdot Q(\mathsf{ct}^r)$.
\end{enumerate}
The total number of non-scalar homomorphic multiplications is
\[
O\!\left(\sqrt{\frac{|S_A(B)|}{r}}\right),
\]
compared to $O(\sqrt{|S_A(B)|})$ for direct Paterson-Stockmeyer evaluation of $P_A$.
The asymptotic speedup factor is $\sqrt{r}$.
\end{theorem}

\subsection{Universality for Large-Prime BGV Bootstrapping}

\begin{corollary}[Universality]\label{cor:universal}
Let $p$ be a prime used in BGV bootstrapping with sparse secret key of Hamming weight $h$.
Let $B = \lfloor k\sqrt{h\phi(m) \cdot 2^{\omega(m)} / (12m)} + 0.5 \rfloor$ be the Ma--Huang--Wang--Wang support bound.

\begin{enumerate}
\item If $p \equiv 1 \pmod{4}$ and $p > (2B+1)^2$, then the order-$4$ cleaner is applicable with
\[
|T| = \frac{|S_A(B)| - 1}{4} + 1, \quad \deg(Q) = \frac{|S_A(B)| - 4}{4}, \quad \text{speedup} = 2\times.
\]

\item If $6 \mid p-1$ and $p > (2B+1)^3$, then the order-$6$ cleaner is applicable with
\[
|T| = \frac{2(|S_A(B)| - 1)}{6} + 1, \quad \deg(Q) = \frac{|S_A(B)| - 6}{6}, \quad \text{speedup} = \sqrt{6} \approx 2.45\times.
\]

\item If $8 \mid p-1$ and $p > (2B+1)^4$, then the order-$8$ cleaner is applicable with
\[
|T| = \frac{3(|S_A(B)| - 1)}{8} + 1, \quad \deg(Q) = \frac{|S_A(B)| - 8}{8}, \quad \text{speedup} = 2\sqrt{2} \approx 2.83\times.
\]
\end{enumerate}

In particular, for all practical large-prime BGV parameters with $p \geq 65537$ and standard sparse key ($h = 12$, $B \leq 17$), condition (1) is satisfied since $(2 \cdot 17 + 1)^2 = 1225 < 65537$.
\end{corollary}

\begin{remark}[Comparison with Prior Work]
\begin{itemize}
\item Chen--Han~\cite{CH18}: degree $(e-1)(p-1)+1$, no structural decomposition.
\item GIKV~\cite{GIKV23}: degree reduced via null-polynomial lattice CVP, no Galois structure.
\item Ma et al.~\cite{Ma24}: degree reduced via bounded-support null polynomials, generic evaluation.
\item \textbf{This work}: same degree as Ma et al., but with Galois-covariant structure enabling $\sqrt{r}$-factor evaluation speedup through the decomposition $P_A(X) = c_1 X + X^{r-1} Q(X^r)$.
\end{itemize}
The key distinction is that our optimization targets the \emph{evaluation cost} (circuit complexity) rather than the \emph{polynomial degree} (representation complexity).
These two optimizations are orthogonal and can be composed.
\end{remark}

\begin{remark}[Tight Optimality via Schur's Lemma]
The term count $|T| = (d-1)(|S_A(B)|-1)/r + 1$ is not merely an upper bound but an \emph{exact value}.
This follows from Schur's lemma applied to the $\langle A \rangle$-module structure of the polynomial space:
the cleaner $P_A$ lies in $\mathrm{Hom}_G(V, \chi_1)$ where $G = \langle A \rangle \cong \mathbb{Z}/r\mathbb{Z}$ acts on $V = \mathbb{F}_p[X]/(G_{S_A})$ and $\chi_1$ is the character determined by the covariance relation.
By Schur's lemma, $\dim \mathrm{Hom}_G(V, \chi_1)$ equals the multiplicity of $\chi_1$ in $V$, which is precisely $|T|$.
No Galois-covariant cleaner with fewer non-zero terms exists on the same support.
\end{remark}
